\section{Goals}

The primary goal of this thesis is to develop and evaluate an intelligent agent system that simplifies interaction between users and Slurm-managed HPC clusters. Specifically, the objectives are:

\begin{itemize}
    \item \textbf{Natural Language Interface.} Design and implement a conversational interface that allows users to query cluster status, inspect accounting and resource limits, submit jobs, manage jobs, and troubleshoot errors using natural language instead of memorizing complex command syntax.
    
    \item \textbf{Observer/Operator Control Flow.} Develop a Slurm-specific workflow where read-only observation is separated from state-changing operation, with structured handoffs and confirmation requirements for risky actions.
    
    \item \textbf{Risk-Aware MCP Integration.} Implement Model Context Protocol (MCP) tools that expose Slurm operations through typed interfaces while preserving the risk distinction between monitoring, diagnosis, submission, job control, and administrative actions.
    
    \item \textbf{Benchmark Dataset \& Evaluation System.} Construct a 3{,}135-case benchmark dataset covering eleven Slurm operational categories (read, diagnose, action, bulk, safety, submission, multi-step, account, edge, docs, domain) across five cluster-state scenarios. Design and implement a deterministic evaluation stack that measures tool recall, routing correctness, HITL compliance, and state-transition accuracy --- enabling architecture-level diagnosis of agent failures beyond simple response-quality scoring.

    \item \textbf{Domain-Adapted Fine-Tuning.} Distill commercial-model traces into role-scoped training data; fine-tune an open-weight model (Qwen2.5-14B-Instruct) via QLoRA to enable local deployment that preserves data sovereignty and eliminates per-query API costs; and evaluate the fine-tuned model against the commercial baseline to quantify the performance--cost trade-off.
    
\end{itemize}
